\documentclass[10pt,a4paper]{amsart}
\usepackage[margin=19mm]{geometry}
\usepackage{amsmath,amssymb,mathtools}
\usepackage[scr=rsfs,cal=euler]{mathalfa}
\usepackage{xfrac}
\usepackage[unicode,hidelinks,bookmarks=false]{hyperref}
\newcommand{\ie}{i.e.\ }
\newcommand{\cf}{cf.\ }
\newcommand{\resp}{respectively\ }
\newcommand{\Subsection}{\subsection}
\providecommand{\nobreakdash}{\nobreak\hbox{-}}
\setlength{\emergencystretch}{2em}
\multlinegap0pt
\usepackage{bm}
\usepackage{bbm}
\usepackage{dsfont}
\usepackage{xspace}
\usepackage{cancel}
\usepackage{mathtools}
\usepackage{amsthm}
\makeatletter
\@ifundefined{thm}{\newtheorem{thm}{Thm}}{}
\@ifundefined{lem}{\newtheorem{lem}[thm]{Lemma}}{}
\makeatother
\makeatletter
\newcommand\be{\begin{equation*}}
\expandafter\ifx\csname customthm\endcsname\relax
\newenvironment{customthm}[1]
  {\renewcommand\theinnercustomthm{#1}\innercustomthm}
  {\endinnercustomthm}
\fi
\newcommand\ee{\end{equation*}}
\makeatother
\title[Rigid-flexible values for symplectic embeddings of four-dime]{Rigid-flexible values for symplectic embeddings of four-dimensional ellipsoids into almost-cubes}
\author{Andrew Lee, Cory H. Colbert}
\date{Source: arXiv:2402.16223v3 (CC BY 4.0)}
\begin{document}
\maketitle

\noindent\emph{Source and licence.} Rigid-flexible values for symplectic embeddings of four-dimensional ellipsoids into almost-cubes. Version arXiv:2402.16223v3 (CC BY 4.0), distributed under the Creative Commons Attribution 4.0 International licence, as recorded in the arXiv licence metadata for this exact version and shown on the abstract page https://arxiv.org/abs/2402.16223v3. Full rights notice: licenses/arxiv-2402.16223-CC-BY-4.0.txt.\par\smallskip

\noindent\textit{Editorial scope.} Counted unit: the Lem whose source label is \texttt{sktksupports} in the article (source0.tex lines 839--847, source label \texttt{sktksupports}), delivered together with the article's own complete original proof of it (source0.tex lines 849--870, one \texttt{proof} environment of 22 lines). No mathematics is written, repaired or re-derived: statement and proof are the article's own text line for line. Required context retained uncounted: none. Every definition, display and label of those lines is retained unchanged. Omitted from this package and not redistributed: 1--838, 848--848, 871--1327. Nothing inside a retained excerpt is omitted or altered: every retained body is delivered exactly as the article's lines are written, so no omission receipt of any kind accompanies this package. Citations: the retained excerpts carry no citation command at all, so no bibliography entry of the article is transcribed and no reference is redistributed. Reference adaptation: none was needed and none was made; every reference inside the delivered unit resolves inside it, so no unresolved reference or citation is produced. Printed slips kept literal: no printed word, symbol, hypothesis or number of the counted statement or of its proof was altered. Layout and class: the article's own class is \texttt{amsart} with packages including graphicx, comment, hyperref, biblatex, ulem, tikz, pgfplots, multicol, standalone, xcolor, enumitem, algorithm2e, geometry, caption; the wrapper is the lane's pinned \texttt{amsart} editorial wrapper at 10pt on A4, which restates the theorem-like environments the retained text needs and adds the article's own macro definitions that the retained text needs (\texttt{\textbackslash be}, \texttt{\textbackslash ee}), with extra wrapper packages (bm, bbm, dsfont, xspace, cancel, mathtools). The delivered numbering is the unit's own. Assets: no retained span contains a figure environment, a graphics-inclusion command, a PDF-page inclusion, a table or a TikZ picture; no image file is copied into this package, the package declares no asset dependency and no image file is redistributed with it. Licence: version arXiv:2402.16223v3 of this article is distributed under the Creative Commons Attribution 4.0 International licence, as recorded in the arXiv licence metadata for this exact version and shown on the abstract page https://arxiv.org/abs/2402.16223v3. A full rights notice is delivered at licenses/arxiv-2402.16223-CC-BY-4.0.txt. Characters: every retained excerpt is pure ASCII, so the delivered engine, XeLaTeX with its default Latin Modern text font, reports no missing character.
\begin{lem}\label{sktksupports}
Let $b \in (1, 2).$ We have the following:
\begin{enumerate}
    \item $\mu_b(T_k)(8) > \frac{2}{\sqrt{b}}$ if and only if $b \in I_{2k}.$
    \item $\mu_b(S_{k+1})(8) > \frac{2}{\sqrt{b}}$ if and only if $b \in I_{2k+1}.$
    \item If $b \in I_{2k},$ then $\mu_b(T_k)(8) \ge \mu_b(C)(8)$ for every $C \in \mathcal{E}_K.$
    \item If $b \in I_{2k+1},$ then $\mu_b(S_{k+1})(8) \ge \mu_b(C)(8)$ for every $C \in \mathcal{E}_K.$
\end{enumerate}
\end{lem}

\begin{proof}
To see (1), consider for each $k \ge 1$ the function $f_{T_k}$ defined below:
\begin{eqnarray*}
f_{T_k}(b) & := & \mu_b(T_k)(8)-\frac{2}{\sqrt{b}} = \frac{8k^2+8k+1}{(2k+1)(k+1) + b\cdot k(2k+1)} - \frac{2}{\sqrt{b}} \\
% f_{S_k}(b) & := & \mu_b(S_k)(8)-\frac{2}{\sqrt{b}} = \frac{8k^2-1}{2k^2+k+ b\cdot (2k^2-k)} - \frac{2}{\sqrt{b}}. \\
\end{eqnarray*}

Observe that $f_{T_k}(b) > 0$ if and only if \begin{equation}\label{sqrtb}
(8k^2 + 8k +1)\sqrt{2b} - \sqrt{8}((2k+1)(k + 1) + b(k(2k+1)) > 0.
\end{equation} Consider the polynomial \be P(t) = (-4k^2-2k)t^2 + (8k^2+8k+1)t + (-4k^2-6k-2).\ee Note that if $t = \sqrt{b},$ we obtain the expression in \ref{sqrtb} (scaled by $1/\sqrt{2}$). The roots of $P$ are $(2k+2)/(2k+1)$ and $(2k+1)/2k.$ Since the leading coefficient of $P$ is negative, and the roots of $P$ are distinct, it follows that $P$ is positive if and only if \be\frac{2k+2}{2k+1} < t < \frac{2k+1}{2k}.\ee Equivalently, since $b>1$ by hypothesis, the expression in \ref{sqrtb} is positive if and only if $\left(\frac{2k+1}{2k+1}\right)^2 < b < \left(\frac{2k+1}{2k}\right)^2,$ or equivalently if and only if $b \in I_{2k}.$

Similarly for (2), the analogous function $f_{S_{k+1}}(b):= \mu_b(S_{k+1})(8)-\frac{2}{\sqrt{b}}>0$ if and only if $b \in I_{2k+1}.$

To see that the obstruction from $T_k$ exceeds that of any other $\{S_i\}\cup \{T_i\}_{i\neq k}$ for these values of $b$, take any $b\in I_{2k}$ and consider any other class $C\in\{S_i\}\cup \{T_i\}_{i\neq k}$:

\begin{eqnarray*}
\mu_b(T_k)(8) - \mu_b(C)(8) & = &  \mu_b(T_k)(8) - \mu_b(C)(8) + \frac{2}{\sqrt{b}} - \frac{2}{\sqrt{b}} \\
& = & f_{T_k}(b) - f_C(b).
\end{eqnarray*}

By the preceding argument, $T_k$ is the unique class such that $f_{T_k}(b)>0$ for $b\in I_{2k}$, so $f_{T_k}(b) - f_C(b) > 0.$ Finally, checking that $\mu_b(\langle 3, 1; 1^{\times 7}\rangle)(8)$ is less than volume on $I_{2k}$ amounts to verifying that \be\frac{2}{\sqrt{b}} - \frac{7}{3+b} > 0,\ee which is straightforward. Consequently, $\mu_b(T_k)(8) \ge \mu_b(C)(8)$ for all $C \in \mathcal{E}_K$ on $I_{2k}.$ The argument for $S_{k+1}$ is identical.
\end{proof}

\end{document}
